\documentclass[11pt,a4paper]{article}

% ============================================================
% Packages
% ============================================================
\usepackage[utf8]{inputenc}
\usepackage[T1]{fontenc}
\usepackage{lmodern}
\usepackage[a4paper,left=0.68in,right=0.68in,top=0.70in,bottom=0.85in,includeheadfoot,headheight=14pt,footskip=24pt]{geometry}
\usepackage{microtype}
\usepackage{amsmath,amssymb,mathtools}
\usepackage{booktabs,tabularx,array,adjustbox,xltabular,makecell,ragged2e,etoolbox,needspace,placeins,hyphenat}
\usepackage{xcolor}
\usepackage{xurl}
\usepackage{hyperref}
\usepackage{enumitem}
\usepackage{listings}
\usepackage{caption}
\usepackage{float}
\usepackage{fancyhdr}
\usepackage{lastpage}

\hypersetup{colorlinks=true,linkcolor=blue!45!black,urlcolor=blue!45!black,citecolor=blue!45!black}

% ============================================================
% Robust anti-overflow layout
% ============================================================
\setlength{\parindent}{0pt}
\setlength{\parskip}{4pt}
\raggedbottom
\sloppy
\emergencystretch=4em
\tolerance=9999
\hbadness=9999
\vfuzz=2pt
\hfuzz=1pt
\allowdisplaybreaks

\setlist[itemize]{leftmargin=*,itemsep=2pt,topsep=2pt,parsep=0pt}
\setlist[enumerate]{leftmargin=*,itemsep=2pt,topsep=2pt,parsep=0pt}

\setlength{\tabcolsep}{2.5pt}
\renewcommand{\arraystretch}{1.08}
\newcolumntype{Y}{>{\RaggedRight\arraybackslash\hspace{0pt}}X}
\newcolumntype{L}[1]{>{\RaggedRight\arraybackslash\hspace{0pt}}p{#1}}

\AtBeginEnvironment{xltabular}{\scriptsize\setlength{\tabcolsep}{2.5pt}\renewcommand{\arraystretch}{1.08}}
\AtEndEnvironment{xltabular}{\normalsize}

\lstdefinestyle{code}{
  basicstyle=\ttfamily\footnotesize,
  breaklines=true,
  breakatwhitespace=false,
  columns=fullflexible,
  keepspaces=true,
  showstringspaces=false,
  frame=single,
  framesep=3pt,
  xleftmargin=0pt,
  xrightmargin=0pt,
  linewidth=\linewidth,
  rulecolor=\color{black!20},
  backgroundcolor=\color{black!2},
  aboveskip=4pt,
  belowskip=4pt
}
\lstset{style=code}

\pagestyle{fancy}
\fancyhf{}
\lhead{DEX Import Documentation}
\rhead{\thepage/\pageref{LastPage}}
\renewcommand{\headrulewidth}{0.3pt}

\title{\textbf{DEX Import Documentation}\\
\large Uniswap v3 USDC/WETH 0.05\% Pool --- Raw Data Import Layer\\[4pt]
\normalsize Version 2.0 \quad|\quad May 2026}
\author{Arthur Gallo}
\date{May 2026}

\begin{document}
\maketitle
\thispagestyle{fancy}

\begin{abstract}
This document describes the decentralised-exchange raw-data import layer used
for the AMM thesis project. It records which Python script imports each
dataset, where the output is written, the file format, the column-level data
dictionary, the time coverage, and the role of each dataset in the later
calibration of AMM, LP, trader, arbitrage, liquidity, and gas-cost models.
The documentation is intentionally restricted to the raw import layer.
Cleaning, aggregation, merging with centralised-exchange data, and
construction of model-ready variables should be performed later in
\texttt{scripts/process\_data/DEX}.

\medskip
\noindent\textbf{Important notes on accuracy.}
Every column description, constant value, formula, and pagination strategy in
this document has been verified directly against the source scripts
(\texttt{dex\_utils.py}, \texttt{pool\_addresses.py},
\texttt{fetch\_all\_monthly.py},
\texttt{fetch\_uniswap\_extra\_events.py},
\texttt{fetch\_uniswap\_pool\_timeseries.py},
\texttt{fetch\_uniswap\_pool\_001pct\_timeseries.py},
\texttt{fetch\_uniswap\_pool\_030pct\_timeseries.py},
\texttt{fetch\_uniswap\_pool\_100pct\_timeseries.py},
\texttt{fetch\_uniswap\_positions.py},
\texttt{fetch\_uniswap\_tick\_snapshots.py},
\texttt{run\_dex.py}, \texttt{run\_multitier\_dex.py}).
Where the scripts differ from a naive reading of the Uniswap v3 protocol,
the script behaviour is documented and the discrepancy is flagged.
\end{abstract}

\tableofcontents
\newpage

% ============================================================
\section{Abbreviations and Terminology}
\label{sec:glossary}

\begin{xltabular}{\linewidth}{L{3.0cm}Y}
\toprule
\textbf{Term} & \textbf{Definition} \\
\midrule
AMM & Automated market maker. A protocol in which a deterministic pricing
      function replaces a traditional order book. \\
WETH & Wrapped Ether. An ERC-20 token pegged 1:1 to ETH. In this pool,
      token1 is WETH; columns ending in \texttt{\_eth} refer to WETH
      amounts, interpreted economically as ETH exposure. \\
USDC & USD Coin. A USD-pegged ERC-20 stablecoin. In this pool, token0 is
      USDC. \\
LP & Liquidity provider. An agent who deposits tokens into a Uniswap v3
     pool position and earns fees. \\
DEX & Decentralised exchange. \\
CEX & Centralised exchange (e.g.\ Binance). \\
TVL & Total value locked. The USD value of tokens currently held in the
      pool. \\
EOA & Externally owned account. A non-contract Ethereum address. \\
Tick & A discrete price level in Uniswap v3, indexed by an integer $i$,
      such that the price at tick $i$ satisfies $p(i)=1.0001^{i}$ (token1
      per token0, in raw token units). \\
Tick spacing & The minimum tick-index increment for which a position can be
      opened. For the 0.05\% fee tier the tick spacing is 10. \\
$L$ & Uniswap v3 virtual liquidity. A non-negative integer satisfying
      $L = \sqrt{\Delta x \cdot \Delta y}$ within a price range, where
      $\Delta x$ and $\Delta y$ are the token reserves that would be
      available if the price traversed the full range. \\
$\sqrt{P}$ & The square root of the pool price, stored as a Q64.96
             fixed-point integer (\texttt{sqrtPriceX96}). \\
Q64.96 & A fixed-point number format with 64 integer bits and 96 fractional
         bits. The true price is $p = (\texttt{sqrtPriceX96}/2^{96})^{2}$. \\
Q128.128 & A fixed-point number format with 128 integer and 128 fractional
           bits. Used for \texttt{feeGrowthInside} accumulators. \\
EIP-1559 & Ethereum Improvement Proposal 1559 (activated August 2021),
           which introduced a base fee burned by the network and a priority
           tip paid to validators. \\
The Graph & A decentralised indexing protocol. Data are sourced from The
            Graph's hosted gateway for the Uniswap v3 subgraph. \\
GraphQL & Query language used by The Graph for data retrieval. \\
MEV & Maximal extractable value. Value extracted by reordering, inserting,
      or censoring transactions within a block. \\
\bottomrule
\end{xltabular}

% ============================================================
\section{Scope}
\label{sec:scope}

The DEX import layer downloads data for the Uniswap v3 USDC/WETH 0.05\%
pool on Ethereum mainnet. The pool address used throughout all scripts is
\[
  \texttt{0x88e6a0c2ddd26feeb64f039a2c41296fcb3f5640}.
\]
The token assignment is:
\[
  \text{token0} = \text{USDC}, \qquad \text{token1} = \text{WETH}.
\]
Several output columns use the suffix \texttt{\_eth} for readability; these
refer to WETH amounts, interpreted economically as ETH exposure.

Raw data are imported from The Graph's Uniswap v3 subgraph through GraphQL
queries. The import layer covers:

\begin{itemize}
  \item swap events (Section~\ref{sec:core});
  \item mint events (Section~\ref{sec:core});
  \item burn events (Section~\ref{sec:core});
  \item collect events (Section~\ref{sec:extra});
  \item flash events (Section~\ref{sec:extra});
  \item hourly and daily pool time series (Section~\ref{sec:timeseries});
  \item current pool metadata and ETH/USD bundle (Section~\ref{sec:timeseries});
  \item current position-level state (Section~\ref{sec:positions});
  \item current and inferred month-end tick-liquidity snapshots
        (Section~\ref{sec:ticks}).
\end{itemize}

% ============================================================
\section{Project Paths}
\label{sec:paths}

The intended project structure is:

\begin{lstlisting}
C:\Courses\thesis_AMM\
    scripts\
        import_data\
            run_dex.py
            run_multitier_dex.py
            DEX\
                dex_utils.py
                pool_addresses.py
                fetch_all_monthly.py
                fetch_uniswap_extra_events.py
                fetch_uniswap_pool_timeseries.py
                fetch_uniswap_pool_001pct_timeseries.py
                fetch_uniswap_pool_030pct_timeseries.py
                fetch_uniswap_pool_100pct_timeseries.py
                fetch_uniswap_positions.py
                fetch_uniswap_tick_snapshots.py
    data_raw\
        DEX\
        multitier\
            fee_100\
            fee_3000\
            fee_10000\
    data_processed\
        DEX\
    docs\
\end{lstlisting}

The main raw-output folder is:
\begin{lstlisting}
C:\Courses\thesis_AMM\data_raw\DEX
\end{lstlisting}

\noindent\textbf{Path resolution.} All scripts resolve the output directory
relative to the script file using
\texttt{Path(\_\_file\_\_).resolve().parents[3] / "data\_raw" / "DEX"}.
This means the scripts must reside exactly three directory levels below the
project root (i.e.\ in
\texttt{scripts\textbackslash import\_data\textbackslash DEX\textbackslash}).
If the scripts are moved, the \texttt{.parents[3]} offset must be updated
accordingly in every affected script, or the output path overridden manually.

\noindent\textbf{Shared utilities.} All download scripts in
\texttt{scripts\textbackslash import\_data\textbackslash DEX\textbackslash}
import \texttt{dex\_utils.py}, which must be placed in the same directory
and must be importable from there (i.e.\ that directory must be on
\texttt{sys.path}, or the scripts must be run from there).
Pool addresses for all four fee tiers are centralised in
\texttt{pool\_addresses.py}, which is imported by the multi-tier fetch
scripts.

% ============================================================
\section{Common Configuration (\texttt{dex\_utils.py})}
\label{sec:config}

All shared constants, the HTTP helper, type-safe coercions, and the
month-range builder live in \texttt{dex\_utils.py}.

\begin{xltabular}{\linewidth}{L{4.3cm}Y}
\toprule
\textbf{Object} & \textbf{Value / description} \\
\midrule
Pool & Uniswap v3 USDC/WETH 0.05\% pool on Ethereum mainnet. \\
Pool address (\texttt{POOL}) & \texttt{0x88e6a0c2ddd26feeb64f039a2c41296fcb3f5640} \\
Fee tier & \texttt{500}, corresponding to $0.05\%$ or 5 basis points. The
           fee rate as a decimal is \texttt{feeTier / 1{,}000{,}000 = 0.0005}. \\
Subgraph ID (\texttt{SUBGRAPH}) &
  \texttt{5zvR82QoaXYFyDEKLZ9t6v9adgnptxYpKpSbxtgVENFV} \\
Subgraph endpoint (\texttt{URL}) &
  \texttt{https://gateway.thegraph.com/api/\{API\_KEY\}/subgraphs/id/\{SUBGRAPH\}} \\
API key & Read at module load time from the environment variable
          \texttt{THEGRAPH\_API\_KEY} via \texttt{\_load\_api\_key()}.
          If the variable is absent or empty, a \texttt{RuntimeError} is
          raised immediately with a human-readable message. API keys must
          not be hard-coded in scripts or committed to version control. \\
Start timestamp (\texttt{START\_TS}) &
  \texttt{1620172800} --- 2021-05-05 00:00:00 UTC.
  This is the first full day after the pool's inception date; it is not
  necessarily the exact deployment block. \\
End timestamp (\texttt{END\_TS}) &
  \texttt{1777593599} --- 2026-04-30 23:59:59 UTC. \\
Batch size (\texttt{BATCH}) &
  \texttt{1000} --- the maximum number of entities per GraphQL page
  (The Graph hard limit). \\
Sleep between queries (\texttt{SLEEP}) &
  \texttt{0.35} seconds. Applied after each successful page. \\
Maximum retries (\texttt{MAX\_RETRIES}) & \texttt{10} (updated from 5 in
  earlier documentation). \\
Base retry wait (\texttt{RETRY\_WAIT}) & \texttt{15} seconds (updated from
  10 seconds). Actual wait uses \emph{exponential backoff}:
  $w_k = \min(15 \cdot 2^{k-1},\, 120)$ seconds on attempt $k$.
  Retry is triggered for any exception caught by the \texttt{run\_query}
  function (HTTP errors, timeout, GraphQL-level error objects in the
  response, and network failures). \\
Time format in CSVs &
  Timestamps are converted with
  \texttt{pd.to\_datetime(..., unit="s", utc=True)}. \\
Main file format &
  CSV files with headers, comma separators, one observation per row,
  UTF-8 encoding (pandas default on all platforms since version 1.0). \\
\bottomrule
\end{xltabular}

\subsection{Type-Safe Coercions}

\texttt{dex\_utils.py} exposes two coercion helpers used throughout the
flatten functions:

\begin{itemize}
  \item \texttt{safe\_int(x)} --- returns \texttt{None} if \texttt{x} is
        \texttt{None} or the empty string, otherwise \texttt{int(x)}.
  \item \texttt{safe\_float(x)} --- returns \texttt{None} if \texttt{x} is
        \texttt{None} or the empty string, otherwise \texttt{float(x)}.
\end{itemize}

The subgraph returns large integers (liquidity, sqrtPrice, Q128.128 fee
accumulators) as strings to avoid JavaScript 64-bit integer overflow. These
fields are therefore kept as strings in the raw CSV and must be cast to
arbitrary-precision integers during processing.

\subsection{Month-Range Builder}

\texttt{build\_month\_ranges(start\_ts, end\_ts)} returns a list of
\texttt{(label, chunk\_start, chunk\_end)} triples, one per calendar month
within $[\texttt{START\_TS}, \texttt{END\_TS}]$. The label format is
\texttt{YYYY\_MM}. The first and last chunks are clipped to the exact
\texttt{START\_TS} and \texttt{END\_TS} boundaries, so they may be
shorter than a full month.

% ============================================================
\section{Script-to-Output Map}
\label{sec:map}

\begin{xltabular}{\linewidth}{L{4.5cm}L{4.0cm}Y}
\toprule
\textbf{Script} & \textbf{Output files} & \textbf{Purpose} \\
\midrule
\texttt{fetch\_all\_monthly.py}
  & \texttt{swaps\_YYYY\_MM.csv} \newline
    \texttt{mints\_YYYY\_MM.csv} \newline
    \texttt{burns\_YYYY\_MM.csv}
  & Downloads core pool events month by month: swaps for trading flow,
    mints for LP liquidity additions, and burns for LP liquidity removals.
    Uses timestamp-cursor pagination. Partial saves every 50{,}000 rows.
    Existing month files are skipped unless \texttt{OVERWRITE = True}. \\
\addlinespace
\texttt{fetch\_uniswap\_extra\_events.py}
  & \texttt{collects\_YYYY\_MM.csv} \newline
    \texttt{flashes\_YYYY\_MM.csv}
  & Downloads additional pool events: fee collections by LPs and
    flash-loan events. Uses \texttt{id\_gt} cursor pagination, which
    is robust to timestamp ties. \\
\addlinespace
\texttt{fetch\_uniswap\_pool\_timeseries.py}
  & \texttt{pool\_hour\_data.csv} \newline
    \texttt{pool\_day\_data.csv} \newline
    \texttt{pool\_metadata\_current.csv} \newline
    \texttt{bundle\_current.csv}
  & Downloads hourly/daily pool aggregates, current pool metadata, and
    the current ETH/USD bundle. Uses \texttt{id\_gt} cursor pagination
    for time series. No overwrite guard; always re-fetches. \\
\addlinespace
\texttt{fetch\_uniswap\_positions.py}
  & \texttt{positions\_current.csv}
  & Downloads current/lifetime position-level state. Not a historical
    panel. Uses \texttt{id\_gt} cursor pagination. \\
\addlinespace
\texttt{fetch\_uniswap\_tick\_snapshots.py}
  & \texttt{tick\_snapshots/ticks\_current.csv} \newline
    \texttt{tick\_snapshots/pool\_state\_current.csv} \newline
    \texttt{tick\_snapshots/ticks\_YYYY\_MM\_block\_N.csv} \newline
    \texttt{tick\_snapshots/pool\_state\_YYYY\_MM\_block\_N.csv} \newline
    \texttt{tick\_snapshots/month\_end\_blocks\_from\_existing\_files.csv}
  & Downloads current and inferred month-end tick snapshots. Month-end
    blocks are inferred as the maximum \texttt{block\_number} across all
    five event-type CSV files for each month. \\
\addlinespace
\texttt{pool\_addresses.py}
  & (no output)
  & Central registry of all four USDC/WETH pool addresses
    (\texttt{POOL\_001}, \texttt{POOL\_005}, \texttt{POOL\_030},
    \texttt{POOL\_100}) and the structured \texttt{MULTITIER\_POOLS} list.
    Imported by the multi-tier fetch scripts; not a standalone executable. \\
\addlinespace
\texttt{fetch\_uniswap\_pool\_001pct\_timeseries.py}
  & \texttt{multitier/fee\_100/pool\_hour\_data.csv} \newline
    \texttt{multitier/fee\_100/pool\_day\_data.csv} \newline
    \texttt{multitier/fee\_100/pool\_metadata\_current.csv}
  & Downloads hourly and daily pool time series for the USDC/WETH 0.01\%
    pool (\texttt{0xe0554a\ldots}). Output schema is identical to
    \texttt{fetch\_uniswap\_pool\_timeseries.py}. Uses
    \texttt{id\_gt} cursor pagination. \\
\addlinespace
\texttt{fetch\_uniswap\_pool\_030pct\_timeseries.py}
  & \texttt{multitier/fee\_3000/pool\_hour\_data.csv} \newline
    \texttt{multitier/fee\_3000/pool\_day\_data.csv} \newline
    \texttt{multitier/fee\_3000/pool\_metadata\_current.csv}
  & Downloads hourly and daily pool time series for the USDC/WETH 0.30\%
    pool (\texttt{0x8ad599\ldots}). \\
\addlinespace
\texttt{fetch\_uniswap\_pool\_100pct\_timeseries.py}
  & \texttt{multitier/fee\_10000/pool\_hour\_data.csv} \newline
    \texttt{multitier/fee\_10000/pool\_day\_data.csv} \newline
    \texttt{multitier/fee\_10000/pool\_metadata\_current.csv}
  & Downloads hourly and daily pool time series for the USDC/WETH 1.00\%
    pool (\texttt{0x7bea39\ldots}). \\
\addlinespace
\texttt{run\_dex.py}
  & (orchestrates the scripts above)
  & Recommended entry point for the main 0.05\% pool import.
    Runs pool time series, monthly events, extra events, positions, and
    (optionally) tick snapshots in the correct order.
    Accepts \texttt{--skip-ticks} to defer tick snapshots. Sets
    \texttt{PYTHONIOENCODING=utf-8} automatically. \\
\addlinespace
\texttt{run\_multitier\_dex.py}
  & (orchestrates multi-tier scripts)
  & Recommended entry point for the multi-tier extension. Downloads
    time-series data for the 0.01\%, 0.30\%, and 1.00\% pools.
    Accepts \texttt{--001}, \texttt{--030}, \texttt{--100} to run
    a single pool. Sets \texttt{PYTHONIOENCODING=utf-8} automatically. \\
\bottomrule
\end{xltabular}

% ============================================================
\section{Recommended Execution Order}
\label{sec:exec}

\begin{enumerate}
  \item Set the API key (Windows, run once per machine):
\begin{lstlisting}
setx THEGRAPH_API_KEY "your_api_key_here"
\end{lstlisting}
  Restart the terminal after running \texttt{setx}.
  Alternatively, create a \texttt{.env} file at the project root:
\begin{lstlisting}
THEGRAPH_API_KEY=your_api_key_here
\end{lstlisting}
  Verify with \texttt{echo \%THEGRAPH\_API\_KEY\%} (or
  \texttt{echo \$env:THEGRAPH\_API\_KEY} in PowerShell).

  \item Install dependencies from the project root:
\begin{lstlisting}
pip install -r requirements.txt
\end{lstlisting}

  \item Run the main 0.05\% pool import:
\begin{lstlisting}
python scripts\import_data\run_dex.py --skip-ticks
\end{lstlisting}
  This runs pool time series, monthly events (swaps, mints, burns),
  extra events (collects, flashes), and LP positions in the correct
  order. \texttt{--skip-ticks} omits tick snapshots, which are the
  slowest step and are not required for most downstream processing.
  Runtime: 1--3 hours (API rate-limited).

  \item (Optional) Run tick snapshots after step~3 completes:
\begin{lstlisting}
python scripts\import_data\DEX\fetch_uniswap_tick_snapshots.py
\end{lstlisting}
  Must follow step~3. The script scans all five monthly event-type
  files (\texttt{swaps}, \texttt{mints}, \texttt{burns},
  \texttt{collects}, \texttt{flashes}) to infer one representative
  block per month.

  \item Run the multi-tier extension (0.01\%, 0.30\%, 1.00\% pools):
\begin{lstlisting}
python scripts\import_data\run_multitier_dex.py
\end{lstlisting}
  Downloads hourly and daily time series for the three additional fee
  tiers into \texttt{data\_raw\textbackslash multitier\textbackslash}.
  Steps~3 and~5 are independent and can run in parallel.
  To download a single fee tier, use \texttt{--001}, \texttt{--030},
  or \texttt{--100}.
\end{enumerate}

\noindent\textbf{Windows encoding note.} Both runner scripts
(\texttt{run\_dex.py} and \texttt{run\_multitier\_dex.py}) set
\texttt{PYTHONIOENCODING=utf-8} automatically before launching child
processes. No manual configuration is required.

% ============================================================
\section{Core Monthly Events: \texttt{fetch\_all\_monthly.py}}
\label{sec:core}

\subsection{Purpose and Pagination Strategy}

This script downloads all swap, mint, and burn events for the pool over
monthly windows from 2021-05 to 2026-04.

\noindent\textbf{Pagination.} It uses a timestamp-cursor strategy inside
each monthly window:
\[
  \texttt{timestamp\_gt} = \texttt{lastTs}, \quad
  \texttt{timestamp\_lte} = \texttt{month\_end},
\]
sorted ascending by \texttt{timestamp}. After each page, \texttt{lastTs}
is updated to the timestamp of the last row in the page:
\[
  \texttt{lastTs} \leftarrow \texttt{int(batch[-1]["timestamp"])}.
\]
The loop terminates when the page returns fewer than \texttt{BATCH} = 1000
rows, or an empty batch.

\noindent\textbf{Timestamp-cursor risk.} If more than 1000 events share
exactly the same \texttt{timestamp}, this pagination strategy can silently
skip rows. The extra-event script avoids this by using \texttt{id\_gt}
(Section~\ref{sec:extra}). For the swap, mint, and burn files, users should
cross-check aggregate row counts against pool-level \texttt{txCount} totals
from \texttt{pool\_day\_data.csv} as a consistency check.

\noindent\textbf{Progress and recovery.} The script logs progress every
\texttt{LOG\_EVERY = 10{,}000} rows and saves a partial CSV every
\texttt{SAVE\_EVERY = 50{,}000} rows to
\texttt{<output\_file>.partial}. On completion the partial file is deleted.
Months whose three output files already exist and are non-empty are skipped
unless \texttt{OVERWRITE = True}.

\subsection{Outputs}

\begin{lstlisting}
C:\Courses\thesis_AMM\data_raw\DEX\swaps_YYYY_MM.csv
C:\Courses\thesis_AMM\data_raw\DEX\mints_YYYY_MM.csv
C:\Courses\thesis_AMM\data_raw\DEX\burns_YYYY_MM.csv
\end{lstlisting}

\subsection{Price and Fee Conventions in the Swap Flattener}

Before the column dictionary, three conventions encoded in
\texttt{flatten\_swaps} are documented here to prevent misuse.

\paragraph{Fee rate.}
The fee is computed as:
\[
  \texttt{fee\_usd} = \texttt{amountUSD} \times \frac{\texttt{feeTier}}{1{,}000{,}000}
  = \texttt{amountUSD} \times 0.0005.
\]
\texttt{feeTier} is cast to \texttt{int} from the subgraph string before
division. The result is an approximation because (a) the fee is collected
from the \emph{input} token before the swap, so the true fee is applied to
the input amount rather than the USD-equivalent of the entire swap; and
(b) \texttt{amountUSD} is computed by the subgraph using its internal
ETH/USD bundle price at indexing time.

\paragraph{Trade direction.}
\[
  \texttt{buy\_eth} = (\texttt{amount0} > 0).
\]
In this pool, token0 is USDC and token1 is WETH. A positive \texttt{amount0}
means the pool receives USDC from the trader, who therefore receives WETH ---
i.e., the trader is buying ETH/WETH with USDC.

\paragraph{Sqrt price and tick.}
The post-swap Q64.96 square-root price \texttt{sqrtPriceX96} and the
post-swap tick \texttt{tick} are stored as returned by the subgraph
(strings). The true float price (token1 per token0, in raw token units) is:
\[
  p_{\mathrm{raw}} = \left(\frac{\texttt{sqrtPriceX96}}{2^{96}}\right)^{2}.
\]
Because token0 has 6 decimals (USDC) and token1 has 18 decimals (WETH), the
human-readable USDC-per-WETH price is:
\[
  p_{\mathrm{USDC/WETH}} = p_{\mathrm{raw}} \times \frac{10^{18}}{10^{6}}
  = p_{\mathrm{raw}} \times 10^{12}.
\]
The subgraph also provides \texttt{token0Price} and \texttt{token1Price}
fields, which should be verified against the formula above on a sample of
observations before use, since the subgraph price convention can be inverted.

\subsection{Swap File Format: \texttt{swaps\_YYYY\_MM.csv}}

Each row is one Uniswap v3 swap event.

\begin{xltabular}{\linewidth}{L{4.5cm}L{3.0cm}Y}
\toprule
\textbf{Column} & \textbf{Format} & \textbf{Description} \\
\midrule
\texttt{swap\_id} & string & Subgraph swap entity ID (from \texttt{r["id"]}). \\
\texttt{tx\_hash} & string & Ethereum transaction hash (from
  \texttt{transaction.id}). \\
\texttt{block\_number} & integer & Block number (from
  \texttt{int(transaction.blockNumber)}). \\
\texttt{log\_index} & string/integer & Event log index inside the
  transaction, passed through as-is from the subgraph. \\
\texttt{timestamp} & UTC datetime & \texttt{pd.to\_datetime(int(timestamp),
  unit="s", utc=True)}. \\
\texttt{origin} & address & EOA that originated the transaction
  (\texttt{tx.origin} semantics). Useful for identifying MEV bots and
  router contracts. \\
\texttt{sender} & address & \texttt{msg.sender} at the pool contract level;
  typically the router or aggregator. \\
\texttt{recipient} & address & Address receiving the output tokens. \\
\texttt{amount0\_usdc} & float & Signed token0 (USDC) amount:
  negative means the pool pays out USDC; positive means the pool receives
  USDC. Cast from \texttt{float(amount0)}. \\
\texttt{amount1\_eth} & float & Signed token1 (WETH) amount. Opposite sign
  to \texttt{amount0\_usdc} in a standard swap. \\
\texttt{amount\_usd} & float & Swap value in USD reported by the subgraph
  (\texttt{amountUSD}). Computed using the subgraph's internal ETH/USD
  bundle price at indexing time. \\
\texttt{fee\_usd} & float & $\texttt{amount\_usd} \times 0.0005$.
  Approximate; see Section~\ref{sec:core}.3. \\
\texttt{buy\_eth} & boolean & \texttt{True} when
  \texttt{amount0} $>$ 0, i.e.\ when the trader buys WETH with USDC. \\
\texttt{sqrt\_price\_x96} & string & Post-swap Q64.96 square-root price
  as returned by the subgraph. Kept as string to avoid integer overflow.
  Convert with $p = (\texttt{val}/2^{96})^{2}$. \\
\texttt{tick\_at\_swap} & string & Pool tick after the swap, as returned
  by the subgraph field \texttt{tick} on the swap entity. \\
\texttt{gas\_used} & integer & Gas used by the full Ethereum transaction
  (\texttt{int(transaction.gasUsed)}). Transaction-level; not uniquely
  attributable to this event if the transaction contains multiple events. \\
\texttt{gas\_price\_wei} & integer & Gas price in wei
  (\texttt{int(transaction.gasPrice)}). After EIP-1559 (August 2021, within
  this sample), this reflects the effective gas price paid (base fee plus
  priority tip), as indexed by the subgraph. \\
\texttt{gas\_cost\_eth} & float & $\texttt{gas\_used} \times
  \texttt{gas\_price\_wei} / 10^{18}$. Computed in the flatten function. \\
\texttt{pool\_liquidity} & string & Active pool liquidity at or after the
  event as indexed by the subgraph. Kept as string (large integer). \\
\texttt{pool\_sqrt\_price} & string & Pool square-root price
  (\texttt{pool.sqrtPrice}). String. \\
\texttt{pool\_tick} & string/integer & Pool tick (\texttt{pool.tick}),
  passed through as-is. \\
\texttt{pool\_token0\_price} & float & Subgraph \texttt{pool.token0Price}.
  Convention must be verified: cross-check a sample against the tick-derived
  price before use. \\
\texttt{pool\_token1\_price} & float & Subgraph \texttt{pool.token1Price}. \\
\texttt{pool\_fee\_tier} & string & Pool fee tier (\texttt{pool.feeTier}),
  passed through as-is; expected value \texttt{"500"}. \\
\texttt{pool\_volume\_usd} & float & Cumulative pool volume in USD
  (\texttt{pool.volumeUSD}). \\
\texttt{pool\_fees\_usd} & float & Cumulative pool fees in USD
  (\texttt{pool.feesUSD}). \\
\texttt{pool\_tx\_count} & integer & Cumulative pool transaction count. \\
\texttt{pool\_tvl\_usd} & float & Total value locked in USD. \\
\texttt{pool\_tvl\_token0\_usdc} & float & TVL in token0 (USDC)
  (\texttt{pool.totalValueLockedToken0}). \\
\texttt{pool\_tvl\_token1\_eth} & float & TVL in token1 (WETH)
  (\texttt{pool.totalValueLockedToken1}). \\
\texttt{token0\_symbol} & string & Expected: \texttt{USDC}. \\
\texttt{token0\_decimals} & string/integer & Expected: \texttt{6}. \\
\texttt{token0\_derived\_eth} & float & Token0 value in ETH
  (\texttt{token0.derivedETH}). \\
\texttt{token1\_symbol} & string & Expected: \texttt{WETH}. \\
\texttt{token1\_decimals} & string/integer & Expected: \texttt{18}. \\
\texttt{token1\_derived\_eth} & float & Token1 value in ETH
  (\texttt{token1.derivedETH}). Expected $\approx 1.0$ since
  token1 is WETH. \\
\bottomrule
\end{xltabular}

\subsection{Mint and Burn File Format}

Each row is one liquidity-addition event (mints) or one liquidity-removal
event (burns). Both files share the same schema, produced by
\texttt{flatten\_mints\_burns}.

\begin{xltabular}{\linewidth}{L{4.5cm}L{3.0cm}Y}
\toprule
\textbf{Column} & \textbf{Format} & \textbf{Description} \\
\midrule
\texttt{id} & string & Subgraph mint or burn entity ID. \\
\texttt{tx\_hash} & string & Ethereum transaction hash. \\
\texttt{block\_number} & integer & Ethereum block number. \\
\texttt{log\_index} & string/integer & Event log index, passed through
  as-is. \\
\texttt{timestamp} & UTC datetime & Event timestamp. \\
\texttt{origin} & address & EOA originating the transaction. \\
\texttt{owner} & address & LP position owner associated with the event. \\
\texttt{liquidity\_units} & float & Amount of Uniswap v3 liquidity
  $L$ minted or burned (\texttt{float(amount)}). For mints, this is the
  liquidity added; for burns, the liquidity removed. \\
\texttt{amount0\_usdc} & float & Token0 (USDC) amount added or removed. \\
\texttt{amount1\_eth} & float & Token1 (WETH) amount added or removed. \\
\texttt{amount\_usd} & float & USD value of the event as reported by the
  subgraph. \\
\texttt{tick\_lower} & string/integer & Lower tick boundary of the LP
  range, passed through as-is. Should be a multiple of tick spacing 10. \\
\texttt{tick\_upper} & string/integer & Upper tick boundary, passed through
  as-is. Should satisfy \texttt{tick\_lower} $<$ \texttt{tick\_upper}. \\
\texttt{gas\_used} & integer & Gas used by the full transaction. \\
\texttt{gas\_price\_wei} & integer & Gas price in wei. \\
\texttt{gas\_cost\_eth} & float & $\texttt{gas\_used} \times
  \texttt{gas\_price\_wei} / 10^{18}$. \\
\texttt{pool\_liquidity} & string & Active pool liquidity after the event. \\
\texttt{pool\_sqrt\_price} & string & Pool square-root price. \\
\texttt{pool\_tick} & string/integer & Pool tick. \\
\texttt{pool\_token0\_price} & float & Subgraph token0 price. \\
\texttt{pool\_token1\_price} & float & Subgraph token1 price. \\
\texttt{pool\_fee\_tier} & string & Pool fee tier. \\
\texttt{pool\_tvl\_usd} & float & Pool TVL in USD. \\
\texttt{token0\_symbol} & string & Expected: \texttt{USDC}. \\
\texttt{token0\_decimals} & string/integer & Expected: \texttt{6}. \\
\texttt{token1\_symbol} & string & Expected: \texttt{WETH}. \\
\texttt{token1\_decimals} & string/integer & Expected: \texttt{18}. \\
\bottomrule
\end{xltabular}

% ============================================================
\section{Extra Monthly Events: \texttt{fetch\_uniswap\_extra\_events.py}}
\label{sec:extra}

\subsection{Purpose and Pagination Strategy}

This script downloads collect and flash events using \texttt{id\_gt} cursor
pagination within each monthly window:
\[
  \texttt{id\_gt} = \texttt{lastId}, \quad
  \texttt{timestamp\_gte} = \texttt{start\_ts}, \quad
  \texttt{timestamp\_lte} = \texttt{end\_ts},
\]
ordered ascending by \texttt{id}. This is robust to timestamp ties.

Output DataFrames are sorted by
(\texttt{timestamp}, \texttt{block\_number}, \texttt{log\_index},
entity ID) before writing to CSV. Partial saves and skip-if-exists logic
mirror the core script (\texttt{SAVE\_EVERY = 50{,}000},
\texttt{OVERWRITE = False}).

\subsection{Outputs}

\begin{lstlisting}
C:\Courses\thesis_AMM\data_raw\DEX\collects_YYYY_MM.csv
C:\Courses\thesis_AMM\data_raw\DEX\flashes_YYYY_MM.csv
\end{lstlisting}

\subsection{Collect File Format: \texttt{collects\_YYYY\_MM.csv}}

Each row is one fee-collection event by an LP position. The flattener uses
\texttt{safe\_int} and \texttt{safe\_float} throughout, so missing fields
produce \texttt{None} rather than errors.

\begin{xltabular}{\linewidth}{L{4.5cm}L{3.0cm}Y}
\toprule
\textbf{Column} & \textbf{Format} & \textbf{Description} \\
\midrule
\texttt{collect\_id} & string & Subgraph collect entity ID. \\
\texttt{tx\_hash} & string & Ethereum transaction hash. \\
\texttt{block\_number} & integer/None & Block number; \texttt{None} if
  absent from subgraph response. \\
\texttt{log\_index} & integer/None & Event log index. \\
\texttt{timestamp} & UTC datetime & Event timestamp. \\
\texttt{owner} & address & LP owner collecting fees. \\
\texttt{amount0\_usdc} & float/None & Collected token0 (USDC) amount. \\
\texttt{amount1\_eth} & float/None & Collected token1 (WETH) amount. \\
\texttt{amount\_usd} & float/None & Collected value in USD as reported
  by the subgraph. \\
\texttt{tick\_lower} & integer/None & Lower tick of the position range. \\
\texttt{tick\_upper} & integer/None & Upper tick of the position range. \\
\texttt{gas\_used} & integer/None & Gas used by the full transaction. \\
\texttt{gas\_price\_wei} & integer/None & Gas price in wei. \\
\texttt{gas\_cost\_eth} & float/None & $\texttt{gas\_used} \times
  \texttt{gas\_price\_wei} / 10^{18}$; \texttt{None} if either gas
  field is \texttt{None}. \\
\texttt{pool\_id} & string & Pool address. \\
\texttt{pool\_fee\_tier} & string & Pool fee tier. \\
\bottomrule
\end{xltabular}

\subsection{Flash File Format: \texttt{flashes\_YYYY\_MM.csv}}

Each row is one Uniswap v3 flash event. A flash loan in Uniswap v3 allows
any caller to borrow tokens within a single transaction, provided the
amounts are repaid (plus fees) before the transaction ends. Flash events
are useful for identifying non-LP, non-trader activity and for completeness
of pool event accounting.

\begin{xltabular}{\linewidth}{L{4.5cm}L{3.0cm}Y}
\toprule
\textbf{Column} & \textbf{Format} & \textbf{Description} \\
\midrule
\texttt{flash\_id} & string & Subgraph flash entity ID. \\
\texttt{tx\_hash} & string & Ethereum transaction hash. \\
\texttt{block\_number} & integer/None & Ethereum block number. \\
\texttt{log\_index} & integer/None & Event log index. \\
\texttt{timestamp} & UTC datetime & Event timestamp. \\
\texttt{sender} & address & Flash-event sender (\texttt{msg.sender} at
  pool level). \\
\texttt{recipient} & address & Address receiving the flash-borrowed tokens. \\
\texttt{amount0\_usdc} & float/None & Token0 (USDC) amount borrowed. \\
\texttt{amount1\_eth} & float/None & Token1 (WETH) amount borrowed. \\
\texttt{amount\_usd} & float/None & USD value of the borrowed amounts,
  as reported by the subgraph. This is the borrowed value, not the
  repaid value. \\
\texttt{amount0\_paid\_usdc} & float/None & Token0 repaid
  (\texttt{amount0Paid}). Includes the 0.05\% fee on the borrowed
  amount. \\
\texttt{amount1\_paid\_eth} & float/None & Token1 repaid
  (\texttt{amount1Paid}). \\
\texttt{gas\_used} & integer/None & Gas used by the full transaction. \\
\texttt{gas\_price\_wei} & integer/None & Gas price in wei. \\
\texttt{gas\_cost\_eth} & float/None & $\texttt{gas\_used} \times
  \texttt{gas\_price\_wei} / 10^{18}$; \texttt{None} if either gas
  field is \texttt{None}. \\
\texttt{pool\_id} & string & Pool address. \\
\texttt{pool\_fee\_tier} & string & Pool fee tier. \\
\bottomrule
\end{xltabular}

% ============================================================
\section{Pool Time-Series Data: \texttt{fetch\_uniswap\_pool\_timeseries.py}}
\label{sec:timeseries}

\subsection{Purpose and Pagination Strategy}

This script downloads hourly and daily pool aggregates, the current pool
metadata snapshot, and the current subgraph bundle. Hourly and daily data
use \texttt{id\_gt} cursor pagination with
\texttt{periodStartUnix\_gte}/\texttt{lte} and \texttt{date\_gte}/\texttt{lte}
filters respectively. There is no skip-if-exists guard; the script always
re-fetches all four files on each run.

\subsection{Outputs}

\begin{lstlisting}
C:\Courses\thesis_AMM\data_raw\DEX\pool_hour_data.csv
C:\Courses\thesis_AMM\data_raw\DEX\pool_day_data.csv
C:\Courses\thesis_AMM\data_raw\DEX\pool_metadata_current.csv
C:\Courses\thesis_AMM\data_raw\DEX\bundle_current.csv
\end{lstlisting}

\subsection{Hourly File Format: \texttt{pool\_hour\_data.csv}}

Each row is one hourly pool aggregate. The DataFrame is sorted by
\texttt{period\_start\_unix} before writing.

\begin{xltabular}{\linewidth}{L{4.2cm}L{3.0cm}Y}
\toprule
\textbf{Column} & \textbf{Format} & \textbf{Description} \\
\midrule
\texttt{id} & string & Subgraph \texttt{poolHourData} entity ID. \\
\texttt{period\_start\_unix} & integer & Start of the hourly period in
  UNIX seconds (\texttt{safe\_int(periodStartUnix)}). \\
\texttt{datetime} & UTC datetime & \texttt{pd.to\_datetime(ts, unit="s",
  utc=True)}. \\
\texttt{liquidity} & string & Pool liquidity during the hour. String
  (large integer). \\
\texttt{sqrt\_price} & string & Pool square-root price. String. \\
\texttt{tick} & integer/None & Pool tick (\texttt{safe\_int}). \\
\texttt{token0\_price} & float/None & Subgraph token0 price. Convention
  to be verified (see Section~\ref{sec:core}.3). \\
\texttt{token1\_price} & float/None & Subgraph token1 price. \\
\texttt{tvl\_usd} & float/None & Total value locked in USD. \\
\texttt{volume\_token0} & float/None & Token0 trading volume during the
  hour. \\
\texttt{volume\_token1} & float/None & Token1 trading volume during the
  hour. \\
\texttt{volume\_usd} & float/None & USD trading volume during the hour. \\
\texttt{fees\_usd} & float/None & Fees in USD during the hour. \\
\texttt{tx\_count} & integer/None & Number of transactions. \\
\texttt{open} & float/None & Opening price for the OHLC record. The
  subgraph OHLC series uses the pool's \texttt{token0Price} convention;
  verify the token ordering before use. \\
\texttt{high} & float/None & High price. \\
\texttt{low} & float/None & Low price. \\
\texttt{close} & float/None & Closing price. \\
\bottomrule
\end{xltabular}

\subsection{Daily File Format: \texttt{pool\_day\_data.csv}}

Same economic fields as the hourly file, one row per calendar day.
The DataFrame is sorted by \texttt{date\_unix} before writing.

\begin{xltabular}{\linewidth}{L{4.2cm}L{3.0cm}Y}
\toprule
\textbf{Column} & \textbf{Format} & \textbf{Description} \\
\midrule
\texttt{id} & string & Subgraph \texttt{poolDayData} entity ID. \\
\texttt{date\_unix} & integer & Start of the day in UNIX seconds
  (\texttt{safe\_int(date)}). \\
\texttt{date} & date & \texttt{pd.to\_datetime(ts, unit="s",
  utc=True).date()}. \\
\texttt{liquidity} & string & Pool liquidity. \\
\texttt{sqrt\_price} & string & Pool square-root price. \\
\texttt{tick} & integer/None & Pool tick. \\
\texttt{token0\_price} & float/None & Subgraph token0 price. \\
\texttt{token1\_price} & float/None & Subgraph token1 price. \\
\texttt{tvl\_usd} & float/None & Total value locked in USD. \\
\texttt{volume\_token0} & float/None & Daily token0 volume. \\
\texttt{volume\_token1} & float/None & Daily token1 volume. \\
\texttt{volume\_usd} & float/None & Daily USD volume. \\
\texttt{fees\_usd} & float/None & Daily fees in USD. \\
\texttt{tx\_count} & integer/None & Daily transaction count. \\
\texttt{open} & float/None & Daily opening price (token0Price convention). \\
\texttt{high} & float/None & Daily high price. \\
\texttt{low} & float/None & Daily low price. \\
\texttt{close} & float/None & Daily closing price. \\
\bottomrule
\end{xltabular}

\subsection{Current Pool Metadata: \texttt{pool\_metadata\_current.csv}}

One row; current indexed pool state. Fetched via a non-paginated single
query for \texttt{pool(id: \$pool)}.

\begin{xltabular}{\linewidth}{L{4.5cm}L{3.0cm}Y}
\toprule
\textbf{Column} & \textbf{Format} & \textbf{Description} \\
\midrule
\texttt{pool\_id} & string & Pool address. \\
\texttt{created\_at\_timestamp} & integer/None & Pool creation timestamp
  in UNIX seconds. \\
\texttt{created\_at\_datetime} & UTC datetime/None & Pool creation time. \\
\texttt{created\_at\_block\_number} & integer/None & Block at pool
  creation. \\
\texttt{fee\_tier} & integer/None & Pool fee tier (expected: 500). \\
\texttt{liquidity} & string & Current active liquidity. \\
\texttt{sqrt\_price} & string & Current square-root price. \\
\texttt{tick} & integer/None & Current pool tick. \\
\texttt{token0\_price} & float/None & Current subgraph token0 price. \\
\texttt{token1\_price} & float/None & Current subgraph token1 price. \\
\texttt{volume\_usd} & float/None & Cumulative pool volume in USD. \\
\texttt{fees\_usd} & float/None & Cumulative pool fees in USD. \\
\texttt{tx\_count} & integer/None & Cumulative transaction count. \\
\texttt{tvl\_usd} & float/None & Current TVL in USD. \\
\texttt{tvl\_token0} & float/None & Current TVL in token0
  (\texttt{totalValueLockedToken0}). \\
\texttt{tvl\_token1} & float/None & Current TVL in token1. \\
\texttt{token0\_id} & string & Token0 contract address. \\
\texttt{token0\_symbol} & string & Token0 symbol. \\
\texttt{token0\_name} & string & Token0 name. \\
\texttt{token0\_decimals} & integer/None & Token0 decimals. \\
\texttt{token0\_derived\_eth} & float/None & Token0 value in ETH. \\
\texttt{token1\_id} & string & Token1 contract address. \\
\texttt{token1\_symbol} & string & Token1 symbol. \\
\texttt{token1\_name} & string & Token1 name. \\
\texttt{token1\_decimals} & integer/None & Token1 decimals. \\
\texttt{token1\_derived\_eth} & float/None & Token1 value in ETH
  (expected $\approx$ 1.0). \\
\bottomrule
\end{xltabular}

\subsection{Current Bundle: \texttt{bundle\_current.csv}}

One row; fetched via a non-paginated single query for \texttt{bundle(id: "1")}.

\begin{xltabular}{\linewidth}{L{4.2cm}L{3.0cm}Y}
\toprule
\textbf{Column} & \textbf{Format} & \textbf{Description} \\
\midrule
\texttt{id} & string & Bundle ID; always \texttt{"1"} in the Uniswap v3
  subgraph. \\
\texttt{eth\_price\_usd} & float/None & Current ETH/USD reference price
  from the subgraph bundle (\texttt{ethPriceUSD}). This is the price used
  internally by the subgraph to compute all \texttt{\_usd} fields and should
  be compared to a CEX reference to detect large discrepancies. \\
\texttt{downloaded\_at\_utc} & UTC datetime & \texttt{pd.Timestamp.utcnow()}
  at time of download. \\
\bottomrule
\end{xltabular}

% ============================================================
\section{Current Position Data: \texttt{fetch\_uniswap\_positions.py}}
\label{sec:positions}

\subsection{Purpose and Pagination Strategy}

This script downloads current/lifetime position-level state using
\texttt{id\_gt} cursor pagination. \textbf{This is not a historical LP
panel.} Historical LP activity must be reconstructed from mints, burns, and
collects. The output DataFrame is sorted by \texttt{position\_id} before
writing. There is no skip-if-exists guard.

\subsection{Output}

\begin{lstlisting}
C:\Courses\thesis_AMM\data_raw\DEX\positions_current.csv
\end{lstlisting}

\subsection{File Format: \texttt{positions\_current.csv}}

Each row is one indexed Uniswap v3 position entity.

\begin{xltabular}{\linewidth}{L{4.8cm}L{3.0cm}Y}
\toprule
\textbf{Column} & \textbf{Format} & \textbf{Description} \\
\midrule
\texttt{position\_id} & string & Uniswap v3 position NFT token ID
  (\texttt{r["id"]}). \\
\texttt{owner} & address & Current indexed position owner. Subject to
  subgraph indexing lag; for NFT positions, the owner may have changed
  since last subgraph sync. \\
\texttt{pool\_id} & string & Pool address. \\
\texttt{pool\_fee\_tier} & integer/None & Pool fee tier. \\
\texttt{liquidity} & string & Current liquidity of the position. String
  (large integer). Zero for closed positions. \\
\texttt{deposited\_token0} & float/None & Lifetime token0 deposited into
  the position (\texttt{depositedToken0}). \\
\texttt{deposited\_token1} & float/None & Lifetime token1 deposited. \\
\texttt{withdrawn\_token0} & float/None & Lifetime token0 withdrawn. \\
\texttt{withdrawn\_token1} & float/None & Lifetime token1 withdrawn. \\
\texttt{collected\_fees\_token0} & float/None & Lifetime collected fees
  in token0 (\texttt{collectedFeesToken0}). \\
\texttt{collected\_fees\_token1} & float/None & Lifetime collected fees
  in token1. \\
\texttt{collected\_fees\_usd} & float/None & Lifetime collected fees in
  USD (\texttt{collectedFeesUSD}), valued at historical subgraph prices.
  Not directly comparable across positions opened at different times. \\
\texttt{fee\_growth\_inside0\_last\_x128} & string & Last recorded
  \texttt{feeGrowthInside0LastX128} for token0. Q128.128 fixed-point
  number stored as string. Used with the global fee growth accumulators
  to compute uncollected fees:
  $\text{owed}_0 = L \times (\texttt{feeGrowthInside0} -
  \texttt{feeGrowthInside0LastX128}) / 2^{128}$. \\
\texttt{fee\_growth\_inside1\_last\_x128} & string & Same for token1. \\
\texttt{tick\_lower} & integer/None & Lower tick boundary
  (\texttt{tickLower.tickIdx}). \\
\texttt{tick\_lower\_price0} & float/None & Subgraph
  \texttt{tickLower.price0}. \\
\texttt{tick\_lower\_price1} & float/None & Subgraph
  \texttt{tickLower.price1}. \\
\texttt{tick\_upper} & integer/None & Upper tick boundary. \\
\texttt{tick\_upper\_price0} & float/None & Subgraph
  \texttt{tickUpper.price0}. \\
\texttt{tick\_upper\_price1} & float/None & Subgraph
  \texttt{tickUpper.price1}. \\
\texttt{token0\_id} & string & Token0 contract address. \\
\texttt{token0\_symbol} & string & Token0 symbol. \\
\texttt{token0\_decimals} & integer/None & Token0 decimals. \\
\texttt{token1\_id} & string & Token1 contract address. \\
\texttt{token1\_symbol} & string & Token1 symbol. \\
\texttt{token1\_decimals} & integer/None & Token1 decimals. \\
\texttt{creation\_tx\_hash} & string & Transaction hash associated with
  position creation (\texttt{transaction.id}). This is the transaction
  that first minted the position NFT; subsequent liquidity additions to
  the same position are not reflected here. \\
\texttt{creation\_block\_number} & integer/None & Block number at
  position creation. \\
\texttt{creation\_timestamp} & UTC datetime/None & Position creation
  timestamp; \texttt{None} if \texttt{transaction.timestamp} is absent. \\
\bottomrule
\end{xltabular}

% ============================================================
\section{Tick Snapshots: \texttt{fetch\_uniswap\_tick\_snapshots.py}}
\label{sec:ticks}

\subsection{Purpose and Methodology}

This script imports tick-level liquidity snapshots. It downloads:
\begin{enumerate}
  \item the current pool state and current ticks;
  \item one pool-state and tick snapshot per inferred month-end block.
\end{enumerate}

\noindent\textbf{Tick filter.} The script sets
\texttt{LIQUIDITY\_NET\_ONLY = True}, which filters ticks with
\texttt{liquidityNet\_not = "0"}. This is sufficient for active-liquidity
reconstruction via the cumulative-\texttt{liquidityNet} method (see
Section~\ref{sec:use}), because the active liquidity $L$ changes only at
ticks where \texttt{liquidityNet} $\neq 0$. It is not a complete listing of
every initialised tick.

\noindent\textbf{Tick cursor.} Pagination uses
\texttt{tickIdx\_gt = lastTick}, starting from $-887{,}273$ (just below the
Uniswap v3 minimum tick of $-887{,}272$), ordered ascending by
\texttt{tickIdx}.

\noindent\textbf{Block-number queries.} Historical snapshots use the
GraphQL \texttt{block: \{ number: \$blockNumber \}} argument on both the
\texttt{pool} and \texttt{ticks} queries. This requires the subgraph to
support historical block queries; The Graph's hosted gateway does support
this for the Uniswap v3 subgraph.

\noindent\textbf{Month-end block inference.} The script calls
\texttt{find\_month\_end\_blocks(data\_dir)}, which scans all files
matching \texttt{(swaps|mints|burns|collects|flashes)\_YYYY\_MM.csv} in
\texttt{data\_raw/DEX}, reads the \texttt{block\_number} column, and takes
the maximum across all five event types for each month label. This ensures
all five event files contribute if present.

\subsection{Outputs}

\begin{lstlisting}
C:\Courses\thesis_AMM\data_raw\DEX\tick_snapshots\ticks_current.csv
C:\Courses\thesis_AMM\data_raw\DEX\tick_snapshots\pool_state_current.csv
C:\Courses\thesis_AMM\data_raw\DEX\tick_snapshots\month_end_blocks_from_existing_files.csv
C:\Courses\thesis_AMM\data_raw\DEX\tick_snapshots\ticks_YYYY_MM_block_BLOCKNUMBER.csv
C:\Courses\thesis_AMM\data_raw\DEX\tick_snapshots\pool_state_YYYY_MM_block_BLOCKNUMBER.csv
\end{lstlisting}

\subsection{Block Selection File: \texttt{month\_end\_blocks\_from\_existing\_files.csv}}

\begin{xltabular}{\linewidth}{L{4.2cm}L{3.0cm}Y}
\toprule
\textbf{Column} & \textbf{Format} & \textbf{Description} \\
\midrule
\texttt{month\_label} & string & Month label in \texttt{YYYY\_MM} format. \\
\texttt{block\_number} & integer & Maximum \texttt{block\_number} found
  across all five event-type CSV files for that month. \\
\bottomrule
\end{xltabular}

\noindent\textbf{Important caveat.} This is not the exact calendar
month-end block. It is the maximum block observed across all downloaded
pool-event files for the month. For exact month-end snapshots, use an
independent block-by-timestamp source (e.g.\ \texttt{etherscan.io} block
search or the Ethereum JSON-RPC \texttt{eth\_getBlockByNumber} with a
known UTC timestamp).

\subsection{Pool-State Snapshot Format}

Applies to \texttt{pool\_state\_current.csv} and
\texttt{pool\_state\_YYYY\_MM\_block\_BLOCKNUMBER.csv}.

\begin{xltabular}{\linewidth}{L{4.2cm}L{3.0cm}Y}
\toprule
\textbf{Column} & \textbf{Format} & \textbf{Description} \\
\midrule
\texttt{snapshot\_label} & string & \texttt{"current"} or
  \texttt{"YYYY\_MM"}. \\
\texttt{block\_number} & integer/None & Block number for historical
  snapshots; \texttt{None} for the current snapshot. \\
\texttt{pool\_id} & string & Pool address. \\
\texttt{fee\_tier} & integer/None & Pool fee tier. \\
\texttt{liquidity} & string & Active pool liquidity. \\
\texttt{sqrt\_price} & string & Pool square-root price. \\
\texttt{tick} & integer/None & Current pool tick. \\
\texttt{token0\_price} & float/None & Subgraph token0 price. \\
\texttt{token1\_price} & float/None & Subgraph token1 price. \\
\texttt{tvl\_usd} & float/None & Total value locked in USD. \\
\texttt{tvl\_token0} & float/None & Total token0 locked
  (\texttt{totalValueLockedToken0}). \\
\texttt{tvl\_token1} & float/None & Total token1 locked. \\
\texttt{token0\_id} & string & Token0 contract address. \\
\texttt{token0\_symbol} & string & Token0 symbol. \\
\texttt{token0\_decimals} & integer/None & Token0 decimals. \\
\texttt{token1\_id} & string & Token1 contract address. \\
\texttt{token1\_symbol} & string & Token1 symbol. \\
\texttt{token1\_decimals} & integer/None & Token1 decimals. \\
\bottomrule
\end{xltabular}

\subsection{Tick Snapshot Format}

Applies to \texttt{ticks\_current.csv} and
\texttt{ticks\_YYYY\_MM\_block\_BLOCKNUMBER.csv}.
DataFrames are sorted by \texttt{tick\_idx} before writing.

\begin{xltabular}{\linewidth}{L{4.5cm}L{3.0cm}Y}
\toprule
\textbf{Column} & \textbf{Format} & \textbf{Description} \\
\midrule
\texttt{snapshot\_label} & string & \texttt{"current"} or
  \texttt{"YYYY\_MM"}. \\
\texttt{block\_number} & integer/None & Block number; \texttt{None} for
  current. \\
\texttt{tick\_id} & string & Subgraph tick entity ID. \\
\texttt{pool\_address} & string & Pool address (\texttt{poolAddress}
  field). \\
\texttt{tick\_idx} & integer/None & Tick index $i$ such that
  $p(i) = 1.0001^{i}$. \\
\texttt{liquidity\_gross} & string & Gross liquidity referencing this
  tick. When \texttt{liquidityGross} reaches zero, the tick may be
  uninitialised by the protocol. String (large integer). \\
\texttt{liquidity\_net} & string & Net liquidity change when crossing
  this tick in the upward direction; sign reverses for downward crossing.
  Non-zero for all rows when \texttt{LIQUIDITY\_NET\_ONLY = True}.
  String (signed large integer). \\
\texttt{price0} & float/None & Subgraph \texttt{price0} at this tick. \\
\texttt{price1} & float/None & Subgraph \texttt{price1} at this tick. \\
\texttt{created\_at\_timestamp} & integer/None & Tick entity creation
  timestamp (\texttt{safe\_int}). \\
\texttt{created\_at\_block\_number} & integer/None & Tick entity creation
  block number (\texttt{safe\_int}). \\
\bottomrule
\end{xltabular}

% ============================================================
\section{Data Use by Model Component}
\label{sec:use}

\begin{xltabular}{\linewidth}{L{4.0cm}L{4.0cm}Y}
\toprule
\textbf{Model object} & \textbf{Raw source} & \textbf{Construction} \\
\midrule
DEX swap count $N_t^{\mathrm{DEX}}$ &
  \texttt{swaps\_YYYY\_MM.csv} &
  Count swap rows per interval $t$. \\
DEX volume $Q_t^{\mathrm{DEX}}$ &
  \texttt{swaps\_YYYY\_MM.csv}; \texttt{pool\_hour/day\_data.csv} &
  Sum \texttt{amount\_usd} per interval, or use aggregate
  \texttt{volume\_usd}. \\
Trade size $q_t$ &
  \texttt{swaps\_YYYY\_MM.csv} &
  Use \texttt{amount\_usd}, \texttt{amount0\_usdc}, or
  \texttt{amount1\_eth}. \\
Trade direction &
  \texttt{swaps\_YYYY\_MM.csv} &
  Use \texttt{buy\_eth} (True $=$ buy WETH with USDC) or signs of
  \texttt{amount0\_usdc} / \texttt{amount1\_eth}. \\
AMM price $p_t^{\mathrm{AMM}}$ &
  \texttt{sqrt\_price\_x96}, \texttt{tick\_at\_swap},
  \texttt{pool\_token0/1\_price} &
  From \texttt{sqrtPriceX96}: $p = (\texttt{val}/2^{96})^{2} \times
  10^{12}$. Alternatively, use subgraph price fields after verifying the
  token-ordering convention on a sample of observations. \\
Fee flow $F_t$ &
  \texttt{fee\_usd}; \texttt{fees\_usd};
  \texttt{collects\_YYYY\_MM.csv} &
  Approximate gross swap fees with \texttt{amount\_usd} $\times$ 0.0005;
  compare to pool aggregate \texttt{fees\_usd} and fee-collection events. \\
LP entry and exit &
  \texttt{mints\_YYYY\_MM.csv}; \texttt{burns\_YYYY\_MM.csv} &
  Use \texttt{liquidity\_units}, token amounts, ticks, timestamps, and
  owner addresses. \\
LP fee collection &
  \texttt{collects\_YYYY\_MM.csv} &
  Use collected token0/token1/USD amounts by owner and range. \\
Uncollected LP fees &
  \texttt{positions\_current.csv} &
  Use \texttt{fee\_growth\_inside0/1\_last\_x128} with pool-level
  global fee accumulators: $\text{owed} = L \times (\Delta
  \texttt{feeGrowthInside}) / 2^{128}$. \\
LP current ranges &
  \texttt{positions\_current.csv} &
  Use current indexed position liquidity and tick bounds. Not a full
  historical panel. \\
Active liquidity by tick &
  \texttt{ticks\_current.csv};
  \texttt{ticks\_YYYY\_MM\_block\_N.csv} &
  Sort ticks ascending by \texttt{tick\_idx}. Starting from the
  pool-state \texttt{liquidity} at the snapshot, apply cumulative
  \texttt{liquidity\_net} (converted from string to signed integer) to
  reconstruct active liquidity at any price. Verify invariant:
  $\sum_i \texttt{liquidityNet}_i = 0$. \\
Pool TVL and liquidity &
  \texttt{pool\_hour/day\_data.csv};
  \texttt{pool\_metadata\_current.csv};
  pool-state snapshots &
  Use \texttt{tvl\_usd}, \texttt{liquidity}, and tick-state variables. \\
Gas costs &
  All event files &
  Use \texttt{gas\_used}, \texttt{gas\_price\_wei}, and
  \texttt{gas\_cost\_eth}. Convert to USD using the ETH/USD bundle or a
  CEX reference. Note that gas fields are transaction-level; if one
  transaction emits multiple events, the gas cost is not uniquely
  attributable to any single event. \\
Flash activity &
  \texttt{flashes\_YYYY\_MM.csv} &
  Use borrowed amounts, repaid amounts, sender/recipient, and gas cost.
  Compare \texttt{amount0\_usdc}/\texttt{amount1\_eth} with
  \texttt{amount0\_paid\_usdc}/\texttt{amount1\_paid\_eth} to compute the
  implicit flash-loan fee. \\
DEX--CEX deviation $d_t$ &
  DEX event/time-series files plus CEX import files &
  Compute $p_t^{\mathrm{AMM}}$ from \texttt{sqrtPriceX96} (preferred for
  precision) or from hourly/daily \texttt{close}. Merge with CEX reference
  price $p_t^{\mathrm{CEX}}$ by timestamp or nearest previous observation.
  $d_t = (p_t^{\mathrm{AMM}} - p_t^{\mathrm{CEX}}) / p_t^{\mathrm{CEX}}$.
  \\
\bottomrule
\end{xltabular}

% ============================================================
\section{Quality Checks}
\label{sec:qc}

After import, run the following checks before processing:

\begin{enumerate}
  \item \textbf{File completeness.} Confirm that each month from 2021-05
        through 2026-04 has \texttt{swaps}, \texttt{mints}, \texttt{burns},
        \texttt{collects}, and \texttt{flashes} files that are non-empty.

  \item \textbf{Chronological order.} Verify each monthly event file is
        sorted by timestamp, then block number, then log index.

  \item \textbf{No inter-month overlap.} Confirm that for adjacent months
        $m$ and $m+1$, the maximum timestamp in month $m$'s files is
        strictly less than the minimum timestamp in month $m+1$'s files.

  \item \textbf{Duplicate events.} Check uniqueness of \texttt{swap\_id},
        \texttt{id} (mints/burns), \texttt{collect\_id}, and
        \texttt{flash\_id} within each file and across all monthly files
        of the same type.

  \item \textbf{Timestamp bounds.} Confirm each monthly file only contains
        rows inside its calendar month (or within the first/last month's
        clipped window).

  \item \textbf{Pool consistency.} Confirm \texttt{pool\_fee\_tier} $=
        500$ and the pool address matches \texttt{POOL} throughout all files
        where these fields are present.

  \item \textbf{Token consistency.} Confirm token0 is USDC (decimals 6)
        and token1 is WETH (decimals 18) in all files where token metadata
        is present.

  \item \textbf{Swap sign check.} Confirm that \texttt{amount0\_usdc} and
        \texttt{amount1\_eth} have strictly opposite signs for all swap rows
        (one positive, one negative), excluding any edge-case zero amounts.

  \item \textbf{Direction consistency.} Confirm \texttt{buy\_eth} $=
        (\texttt{amount0\_usdc} > 0)$ for all swap rows.

  \item \textbf{Gas sanity.} Check that \texttt{gas\_used},
        \texttt{gas\_price\_wei}, and \texttt{gas\_cost\_eth} are
        non-negative and not systematically missing.

  \item \textbf{Tick bounds.} Confirm \texttt{tick\_lower} $<$
        \texttt{tick\_upper} for all mint, burn, collect, and position
        records where both are non-null. Confirm both are multiples of
        tick spacing 10 (the spacing for a 0.05\% pool).

  \item \textbf{Liquidity units.} Confirm \texttt{liquidity\_units} $> 0$
        for all mint and burn rows (zero-liquidity events should be absent
        under normal subgraph indexing).

  \item \textbf{Aggregate fee consistency.} Compute
        $\sum \texttt{fee\_usd}$ across all swap files and compare to
        cumulative \texttt{fees\_usd} from \texttt{pool\_metadata\_current.csv}
        as a rough upper-bound check.

  \item \textbf{Sqrt-price tick consistency.} For a sample of swap rows,
        compute the tick implied by \texttt{sqrt\_price\_x96} using
        $i = \lfloor \log_{1.0001}(p_{\mathrm{raw}}) \rfloor$ and
        verify it matches \texttt{tick\_at\_swap} within $\pm 1$.

  \item \textbf{Snapshot sanity.} Confirm that month-end snapshot files
        correspond to blocks listed in
        \texttt{month\_end\_blocks\_from\_existing\_files.csv}. Confirm
        $\sum_i \texttt{liquidityNet}_i = 0$ (signed, as integers) for each
        tick snapshot.

  \item \textbf{Bundle price check.} Compare \texttt{eth\_price\_usd}
        in \texttt{bundle\_current.csv} against a contemporaneous CEX
        spot price to detect large subgraph bundle deviations.

  \item \textbf{USDC depeg window.} Flag rows where
        \texttt{timestamp} falls between 2023-03-10 and 2023-03-13 UTC,
        during which USDC briefly traded below \$1.00. All
        \texttt{amount\_usd} and \texttt{tvl\_usd} fields in this window
        should be treated as measured in non-par USDC, not in USD.
\end{enumerate}

% ============================================================
\section{Known Limitations}
\label{sec:limits}

\begin{enumerate}
  \item \textbf{Subgraph dependence.} All files are indexed subgraph data,
        not direct archive-node event logs. The subgraph may have indexing
        gaps, reorgs, or schema-level deviations from on-chain state. For
        high-stakes empirical claims, key aggregates should be compared
        against direct RPC queries or another independent index.

  \item \textbf{Subgraph version.} The subgraph ID
        \texttt{5zvR82QoaXYFyDEKLZ9t6v9adgnptxYpKpSbxtgVENFV} identifies a
        specific deployment of the Uniswap v3 subgraph schema. The schema
        version is not pinned in the scripts. If The Graph re-deploys the
        subgraph under a new ID or schema, field names or types may change.

  \item \textbf{Timestamp-cursor risk in core events.} The core monthly
        script paginates swaps, mints, and burns on timestamp alone.
        If a page boundary falls in the middle of a timestamp tie
        (more than 1000 events at the exact same second), rows can be
        silently skipped. The extra-event script avoids this via
        \texttt{id\_gt}. A more robust version of the core importer would
        use a combined (\texttt{timestamp}, \texttt{id}) cursor.

  \item \textbf{Gas fields are transaction-level.} If one transaction
        triggers multiple pool events (e.g., a swap and a mint in the same
        tx), the \texttt{gas\_used} and \texttt{gas\_price\_wei} values are
        duplicated across all events from that transaction. Event-level gas
        cost is not uniquely identifiable from the raw import.

  \item \textbf{EIP-1559 gas price.} After EIP-1559 (activated block
        12{,}965{,}000, approximately 2021-08-05), the effective gas price
        is $\min(\texttt{maxFeePerGas},\, \texttt{baseFee} +
        \texttt{maxPriorityFeePerGas})$. The subgraph \texttt{gasPrice}
        field is assumed to reflect the effective post-EIP-1559 price,
        but this has not been independently verified against RPC data
        for this version of the subgraph. A verification check on a
        sample of post-1559 transactions is recommended.

  \item \textbf{Current positions are not historical positions.}
        \texttt{positions\_current.csv} gives the current indexed state of
        positions. It may omit positions that were fully withdrawn and
        garbage-collected by the subgraph, and should not be treated as a
        complete historical LP registry. Reconstruct the full LP panel from
        mints, burns, and collects.

  \item \textbf{Month-end snapshots are approximate.} The tick snapshot
        script uses the maximum block number in downloaded event files. This
        is the latest block from pool-active events within the month, which
        may precede the true calendar month-end block if no events occurred
        near month end. For exact calendar month-end snapshots, supply block
        numbers from an independent source.

  \item \textbf{Token price convention must be verified.} Subgraph price
        fields (\texttt{token0Price}, \texttt{token1Price}, \texttt{price0},
        \texttt{price1}) can be inverted relative to the intuitive USDC/WETH
        convention. Before using any price field, verify on a sample of
        observations that the magnitude and direction match an independent
        reference (e.g., Binance ETH/USDT spot price).

  \item \textbf{USD fields use subgraph bundle pricing.}
        All \texttt{\_usd} fields in every file are computed by the subgraph
        using its internal \texttt{ethPriceUSD} bundle at indexing time.
        This may differ from a contemporaneous CEX spot price, introducing
        measurement error in any USD-denominated variable.

  \item \textbf{USDC depeg risk.} USDC briefly depegged from \$1.00 in
        March 2023 (approximately 2023-03-10 to 2023-03-13). During this
        window, all \texttt{\_usd} fields in the subgraph data are measured
        in non-par USDC, not in true USD. Analyses spanning this period
        should apply a USDC/USD correction or exclude the affected dates.

  \item \textbf{MEV and sandwich transactions.} The USDC/WETH 0.05\% pool
        is a primary target for MEV activity, including sandwich attacks and
        arbitrage. The \texttt{origin} and \texttt{sender} fields can be
        used to flag known MEV bot addresses, but not all MEV activity is
        identifiable from on-chain data alone. Trade-size and
        direction-based inference from \texttt{buy\_eth} may be distorted by
        sandwiched transactions that appear as back-to-back opposing trades
        in the same block.

  \item \textbf{WETH vs.\ ETH naming.} The pool trades WETH on-chain.
        Columns ending in \texttt{\_eth} are economic shorthand for WETH
        exposure; no unwrapping is involved.

  \item \textbf{API key handling.} API keys are read from environment
        variables and must not be hard-coded in scripts or committed to
        version control. The \texttt{\_load\_api\_key()} function raises
        immediately if the key is missing.

  \item \textbf{Block reorg risk.} Events indexed close to the chain tip at
        download time may be affected by block reorganisations. Data from
        the last few days before the download date should be treated as
        potentially unstable and re-fetched after confirmation.

  \item \textbf{Path coupling.} All scripts use
        \texttt{Path(\_\_file\_\_).resolve().parents[3]} to locate the
        project root. Moving any script outside
        \texttt{scripts/import\_data/DEX/} will silently write output to an
        incorrect path.
\end{enumerate}

% ============================================================
\section{Recommended Processing Outputs}
\label{sec:proc}

The raw import files are transformed into model-ready datasets by
\texttt{scripts\textbackslash process\_data\textbackslash DEX}.
The actual processed outputs (as of the April 2026 data cut-off) are:

\begin{lstlisting}
data_processed\DEX\dex_swaps.csv             (10,895,792 rows, ~7.6 GB)
data_processed\DEX\dex_pool_hourly.csv       (43,728 rows)
data_processed\DEX\dex_pool_daily.csv        (1,822 rows)
data_processed\DEX\dex_lp_positions.csv      (70,103 rows)
data_processed\DEX\dex_lvr_hourly.csv        (34,008 rows)
data_processed\merged\merged_hourly.csv      (43,728 rows, DEX+CEX joined)
data_processed\calibration\calibration_params.json
data_processed\calibration\calibration_summary.csv
\end{lstlisting}

Core constructed variables:
\[
  N_t^{\mathrm{DEX}},\quad
  Q_t^{\mathrm{DEX}},\quad
  \bar{q}_t^{\mathrm{DEX}},\quad
  p_t^{\mathrm{AMM}},\quad
  \mathrm{fees}_t,\quad
  L_t,\quad
  \mathrm{TVL}_t,\quad
  g_t,\quad
  d_t = \frac{p_t^{\mathrm{AMM}} - p_t^{\mathrm{CEX}}}{p_t^{\mathrm{CEX}}}.
\]

For $p_t^{\mathrm{AMM}}$, the recommended construction is:
\[
  p_t^{\mathrm{AMM}} = \left(\frac{\texttt{sqrt\_price\_x96}}{2^{96}}\right)^{2}
  \times 10^{12} \quad [\text{USDC per WETH}],
\]
where the $10^{12}$ factor corrects for the 12-decimal difference between
USDC (6 decimals) and WETH (18 decimals). Verified to machine precision
(maximum relative error $2.5 \times 10^{-16}$) across all 10.9M swaps.
The LVR rate is computed as $\sigma^{2}_{\mathrm{ann}}/8$ (fraction of TVL
per year).  The hourly/daily OHLC fields use
the subgraph's own price convention and should be verified separately.

\end{document}
